\documentclass[12pt]{article}
% -------------------
% Packages
% -------------------
\usepackage{
	amsmath,			% Math Environments
	amssymb,			% Extended Symbols
	amsthm,			    % Theorem Environments
	cancel,			    % Use Cancels
	enumerate,		    % Enumerate Environments
	graphicx,			% Include Images
	lastpage,			% Reference Lastpage
	multicol,			% Use Multi-columns
	multirow,			% Use Multi-rows
	xcolor,			    % Use Colors
	float,
	geometry,
	quiver
	}

% -------------------
% Header & Footer
% -------------------
\usepackage{fancyhdr}

\newcommand{\assignment}[1]{
	\fancypagestyle{title}{
		%Headers
		\fancyhead[L]{Jake Bridge}
		\fancyhead[C]{MATH 418}
		\fancyhead[R]{HW #1}
		\renewcommand{\headrulewidth}{0.2pt}
		%Footers
		\fancyfoot[L]{}
		\fancyfoot[C]{}
		\fancyfoot[R]{}
		\renewcommand{\footrulewidth}{0.0pt}
	}
	\thispagestyle{title}
}
\fancypagestyle{pages}{
	%Headers
	\fancyhead[L]{}
	\fancyhead[C]{}
	\fancyhead[R]{}
	\renewcommand{\headrulewidth}{0.0pt}
	%Footers
	\fancyfoot[L]{}
	\fancyfoot[C]{}
	\fancyfoot[R]{}
	\renewcommand{\footrulewidth}{0.0pt}
}
\headheight=18pt
\footskip=14pt

\pagestyle{pages}

% ---------------
% Page Layout
% ---------------
\geometry{letterpaper,
bindingoffset=0.2in,
left=1in,
right=1in,
top=1in,
bottom=1in,
footskip=.25in}
% -------------------
% Font
% -------------------
%\usepackage[T1]{fontenc}
%\usepackage{charter}

%\usepackage[T1]{fontenc}
%\usepackage{mathpazo}

%\usepackage[bitstream-charter]{mathdesign}
%\usepackage[T1]{fontenc}


% -------------------
% Commands
% -------------------

% Problem Labels
\newcounter{problem}[section]
\newenvironment{problem}[1][\theproblem]{\refstepcounter{problem}\noindent \textbf{Problem~#1.\quad}}{\vskip\smallskipamount}


\newenvironment{solution}{\noindent \textbf{Solution.\quad}}{\vskip\bigskipamount}


% Special Characters
\newcommand{\N}{\mathbb{N}}
\newcommand{\Z}{\mathbb{Z}}
\newcommand{\Q}{\mathbb{Q}}
\newcommand{\R}{\mathbb{R}}
\newcommand{\C}{\mathbb{C}}

\newcommand{\mc}[1]{\mathcal{#1}}

\newcommand{\ob}[1]{\text{ob}(#1)}
\newcommand{\natt}{\Rightarrow}
\newcommand{\iso}{\cong}
\newcommand{\op}{\text{op}}

\newcommand{\Set}{\textbf{Set}}
\newcommand{\Hom}{\textbf{Hom}}
\newcommand{\Small}{\mathbb{S}}

\newcommand{\eps}{\varepsilon}

% Math Operators

% Special Commands



% DOC
\begin{document}
	\assignment{10}
\newcommand{\G}{\mathbb{G}}
\begin{problem}[104]
	Recall that the category of directed graphs is the category of presheaves $[\G^\op, \Set]$ with $\G = % https://q.uiver.app/#q=WzAsMixbMCwwLCIwIl0sWzEsMCwiMSJdLFswLDEsIlxcdGF1IiwxLHsib2Zmc2V0IjoyfV0sWzAsMSwiXFxzaWdtYSIsMSx7Im9mZnNldCI6LTJ9XV0=
	\begin{tikzcd}[cramped]
		0 & 1
		\arrow["\tau"{description}, shift right=1, from=1-1, to=1-2]
		\arrow["\sigma"{description}, shift left=2, from=1-1, to=1-2]
	\end{tikzcd}$
	
	\begin{enumerate}
		\item Compute $H_0$ and $H_1$ explicitly. Make sure they make sense.
		\item Compute $H_1 \times H_1$. Is it a square?
	\end{enumerate}
\end{problem}


\begin{solution}
	$H_\bullet: \G \to [\G^\op, \Set]$, so $H_0: \G^\op \to \Set$ is given by $H_0(-) = \G(-, 0)$. 
	
	In particular, $H_0(0) = \{1_0\}, H_0(1) = \emptyset$. Thus, because $H_0(0 \to 1)$ has an empty set as its domain (remember to op!),  $H_0(\sigma)$ is the empty function, as is $H_0(\tau)$. This "looks like" a lone vertex unconnected to anything.
	
	Similarly, $H_1(-) = \G(-, 1)$. Then, $H_1(0) = \{\sigma, \tau\}, H_1(1) = \{1_1\}, H_1(\sigma) : \G(1, 1) \to \G(0, 1)$ must be given by prescomposition by $\sigma$, and similar for $\tau$. In particular, because the only morphism in the domain is $1_1$, $H_1(\sigma)(1_1) = \sigma$, and $H_1(\tau)(1_1) = \tau$. This "looks like" a lone edge with a unique source and target.
	
	As we've discussed, $H_1 \times H_1$ as a product (hence limit) can be built from its limit (product) on its components. 
	
	We write $H_1(0) \times H_1(0)$ as a cartesian product as $\{\sigma, \tau\} \times \{\sigma, \tau\} = $\[\{(\sigma, \tau), (\tau, \sigma), (\sigma, \sigma), (\tau, \tau)\}\]
	
	We interpret this as the set of vertices of our graph.
	
	We also write $H_1(1) \times H_1(1) = \{1_1\} \times \{1_1\} = \{(1_1, 1_1)\}$. This is our set of edge (singular).
	
	To find where the edge starts, we take $(H_1(\sigma), H_1(\sigma))(1_1, 1_1) = (\sigma, \sigma)$. To find where the edge ends, we can do the same thing but with $\tau$ and get $(\tau, \tau)$. 
	
	% https://q.uiver.app/#q=WzAsNCxbMCwwLCJcXHNpZ21hXFxzaWdtYSJdLFsyLDAsIlxcc2lnbWFcXHRhdSJdLFswLDEsIlxcdGF1XFxzaWdtYSJdLFsyLDEsIlxcdGF1XFx0YXUiXSxbMCwzXV0=
	\[\begin{tikzcd}
		{\sigma\sigma} && {\sigma\tau} \\
		{\tau\sigma} && {\tau\tau}
		\arrow[from=1-1, to=2-3]
	\end{tikzcd}\]
	
	This is not a square. 
	
	

\end{solution}





\begin{problem}[106]
	Let $F: \mc{A} \to \mc{B}$ and $G: \mc{C} \to \mc{B}$ be functors. The comma category $F/G$ has objects $(a \in \mc{A}, f: Fa \to Gc, c \in \mc{C})$. 
	
	\begin{enumerate}
		\item Finish the definition of the comma category
		\item How are the other comma categories we've defined special cases of this?
		\item Prove that the category of elements of a presheaf is a comma category
	\end{enumerate}
\end{problem}
\begin{solution}
	and another solution. Hallelujah!
\end{solution}


\end{document}
